\par\addvspace{12pt}\Needspace{9\baselineskip}

\section{定积分}

% ITEM: calculus-example-052
\solutionitem{calculus-example-052}{例5.52}
某质点在\(0\le t<1\)时速度为\(2\)，在\(1\le t\le3\)时速度为\(-1\)。求总位移与总路程。再对速度\(v(t)=t\)、\(0\le t\le2\)，分别用四个等长时间段的左、右端点速度估计位移。
\solutionanswer
总位移为\(2\times1-1\times2=0\)，总路程为\(2\times1+1\times2=4\)。第二问每段长\(1/2\)，左端点和为\((0+1/2+1+3/2)/2=3/2\)，右端点和为\((1/2+1+3/2+2)/2=5/2\)。速度递增，故真实位移介于两者之间；由三角形面积可知它等于\(2\)。


% ITEM: calculus-exercise-31-01
\solutionitem{calculus-exercise-31-01}{习题5.31（1）}
用定义求\(\int_a^b(kx+c)\,dx\)，并给出取样和误差关于\(\|P\|\)的界。
\solutionanswer
令\(J=k(b^2-a^2)/2+c(b-a)\)。因\(\sum_i[(x_{i-1}+x_i)/2]\Delta x_i=(b^2-a^2)/2\)，
\[
 |S-J|\le\frac{|k|}2\sum_i(\Delta x_i)^2\le\frac{|k|(b-a)}2\|P\|\to0.
\]
这对所有分割、取样成立，故积分为\(J\)。


% ITEM: calculus-exercise-31-02
\solutionitem{calculus-exercise-31-02}{习题5.31（2）}
设分割\(P_n\)始终保留\([0,1/2]\)这一整段，只将\([1/2,1]\)不断细分。分段数趋于无穷时，\(\|P_n\|\)是否趋零？对\(f(x)=x\)，改变第一段的取样能否产生不消失的差异？
\solutionanswer
不趋零，最大段长始终至少\(1/2\)。保持其余取样不变，第一段分别选\(0\)、\(1/2\)，两和始终相差\((1/2-0)/2=1/4\)。增加分段数不能代替每段都趋短。


% ITEM: calculus-exercise-31-03
\solutionitem{calculus-exercise-31-03}{习题5.31（3）}
对\(f(x)=x^2\)，计算\([0,1]\)上等分分割的左右端点和及其极限。区分这个计算与证明一般定义所需完成的任务。
\solutionanswer
左和为\((n-1)n(2n-1)/(6n^3)\)，右和为\(n(n+1)(2n+1)/(6n^3)\)，均趋\(1/3\)。这只计算了两列特定和；一般定义还要求任意足够细的分割和任意取样都趋同一数。连续函数可积性可以补足这一步。


% ITEM: calculus-example-053
\solutionitem{calculus-example-053}{例5.53}
判断下列函数在\([0,1]\)上是否可积，并说明所用的函数类：\(|x-1/2|\)；在\(x<1/2\)时为\(0\)、其余位置为\(1\)的阶跃函数；在正半轴为\(\sin(1/x)\)、零点补零的函数。比较它们的连续性与可积性。
\solutionanswer
第一函数在闭区间连续，故可积；第二函数有界且仅在\(1/2\)处间断，也可由单调函数可积性得到；第三函数有界，仅在\(0\)处间断，故仍可积。后两者说明连续性是充分条件，并非必要条件。


% ITEM: calculus-exercise-32-01
\solutionitem{calculus-exercise-32-01}{习题5.32（1）}
本题中\(g(0)=0\)，且当\(1/(n+1)<x\le1/n\)、\(n\ge1\)时\(g(x)=1/n\)。
画出上述\(g\)在\([1/4,1]\)上的图像，证明它单调不减，并判断在零点是否连续。
\solutionanswer
在所给区间，\(g(1/4)=1/4\)，\((1/4,1/3]\)上为\(1/3\)，\((1/3,1/2]\)上为\(1/2\)，\((1/2,1]\)上为\(1\)。图像由这些水平段及端点组成，各段左端空心、右端实心，另补\((1/4,1/4)\)。一般地，\(x\)增大时所处段的指标\(n\)不增，故\(g\)不减。又\(0\le g(x)\le2x\)，所以\(g(x)\to0=g(0)\)，在零点连续。
\begin{center}\begin{tikzpicture}[x=5cm,y=2cm]
\draw[->](.18,0)--(1.08,0)node[right]{$x$};\draw[->](.2,0)--(.2,1.2)node[above]{$g$};
\foreach \a/\b/\c in {.25/.333333/.333333,.333333/.5/.5,.5/1/1}{\draw[thick](\a,\c)--(\b,\c);\draw[fill=white](\a,\c)circle(1.4pt);\fill(\b,\c)circle(1.4pt);}
\fill(.25,.25)circle(1.4pt);\node[below]at(.25,0){$1/4$};\node[below]at(.5,0){$1/2$};\node[below]at(1,0){$1$};
\end{tikzpicture}\end{center}


% ITEM: calculus-exercise-32-02
\solutionitem{calculus-exercise-32-02}{习题5.32（2）}
在\([0,1]\)上，令\(h(1/2)=1\)，其余位置\(h(x)=0\)。用取样和定义证明\(h\)可积且积分为零。
\solutionanswer
非零取样只能来自包含\(1/2\)的小段，至多两段，总长不超过\(2\|P\|\)。故\(0\le S(h;P,\xi)\le2\|P\|\to0\)，按定义可积，积分为零。


% ITEM: calculus-exercise-32-03
\solutionitem{calculus-exercise-32-03}{习题5.32（3）}
分别说明“连续”“只有有限个间断点”“单调”“有界”与Riemann可积之间已经建立的蕴含关系，指出哪些条件不能反过来推得。
\solutionanswer
闭区间连续、有限个间断点且有界、闭区间单调，分别都蕴含Riemann可积；可积必有界。有界不能推出可积，如Dirichlet函数；可积不能推出连续，如阶跃函数，不能推出间断点有限，如无穷阶梯函数，也不能推出单调，如\(x(1-x)\)。连续不保证单调，单调不保证连续；可积性的共同要求是取样误差能统一控制。


% ITEM: calculus-example-054
\solutionitem{calculus-example-054}{例5.54}
不使用原函数，利用几何及区间可加性计算\(\int_{-1}^1x\,dx\)、\(\int_{-1}^1|x|\,dx\)和\(\int_0^2|x-1|\,dx\)。解释三个积分的符号与面积意义。
\solutionanswer
第一积分的左右两块三角形面积均为\(1/2\)，符号相反，故积分为零；第二积分把负半轴的一块翻到上方，值为\(1\)；第三积分是以\(x=1\)为分界的两个面积\(1/2\)的三角形，值也为\(1\)。第一个是有向面积的抵消，后两个被积函数非负，积分就是总面积。


% ITEM: calculus-example-055
\solutionitem{calculus-example-055}{例5.55}
不求精确值，利用\(0\le x\le1\)时\(1-x^2\le e^{-x^2}\le1\)，估计\(\int_0^1 e^{-x^2}\,dx\)。再将区间分成两半，分别用每段的函数值上下界估计，比较所得范围。
\solutionanswer
由比较得\(2/3\le I\le1\)，其中\(\int_0^1x^2=1/3\)可由等分和及连续可积性确定。分成两半后，函数递减，端点上下界给出
\[
 \frac{e^{-1/4}+e^{-1}}2\le I\le\frac{1+e^{-1/4}}2.
\]
约为\(0.57334\le I\le0.88940\)。这次上界改善，下界反而较粗；把两种估计合用得\(2/3\le I\le(1+e^{-1/4})/2\)。分段本身并不保证优于已有的非恒定比较函数。


% ITEM: calculus-exercise-33-01
\solutionitem{calculus-exercise-33-01}{习题5.33（1）}
已知\(\int_0^1f=2\)、\(\int_1^3f=-1\)，求\(\int_0^3f\)、\(\int_3^0f\)及\(\int_0^1(3f-2)\)。
\solutionanswer
依次为\(1,-1,4\)，分别由区间可加性、反向积分及线性性得到。


% ITEM: calculus-exercise-33-02
\solutionitem{calculus-exercise-33-02}{习题5.33（2）}
设\(f\)连续且\(1\le f(x)\le3\)，\(x\in[0,2]\)。估计其积分和平均值。若积分等于\(2\)，能否确定\(f\)？
\solutionanswer
积分在\([2,6]\)，平均值在\([1,3]\)。若积分等于\(2\)，则连续非负函数\(f-1\)积分为零。若它某点为正，连续性保证一小段上有正下界，会使积分为正，矛盾。因此\(f\equiv1\)。


% ITEM: calculus-exercise-33-03
\solutionitem{calculus-exercise-33-03}{习题5.33（3）}
设\(f\)在\([a,b]\)上连续。证明\(\int_a^b f^2=0\)当且仅当\(f\equiv0\)。
\solutionanswer
\(f\equiv0\)显然充分。反之，若\(f(c)\ne0\)，由连续性在\(c\)附近的一段正长度区间内\(f^2\ge f(c)^2/4>0\)，故总积分为正，矛盾。端点处取单侧小区间即可。


% ITEM: calculus-exercise-33-04
\solutionitem{calculus-exercise-33-04}{习题5.33（4）}
若\(|f(x)-g(x)|\le10^{-3}\)，\(x\in[-2,3]\)，估计两积分之差。若只知道某一点的误差很小，能否得到同样结论？
\solutionanswer
由积分绝对值估计，差不超过\(5\times10^{-3}\)。单点误差小不足以控制整体：例如\(f(x)-g(x)=M(x+2)\)在\(-2\)处为零，而积分差为\(25M/2\)，可任意大。


% ITEM: calculus-exercise-33-05
\solutionitem{calculus-exercise-33-05}{习题5.33（5）}
构造一个有界的分段常值函数，其平均值不等于任何一点的函数值，说明积分中值定理中连续性的作用。
\solutionanswer
在\([0,1/2)\)取\(0\)，其余位置取\(1\)。积分及平均值都是\(1/2\)，却没有任何点的函数值是\(1/2\)。连续性在积分中值定理中负责保证平均值属于实际值域。
\par
